\subsubsection{Python Utils}
\begin{enumerate}
    \item \textbf{Python BST Interface}
          
          \begin{enumerate}
            \item \textbf{Build BST Index}
                  
                  Python wrapper for building BST index from NumPy arrays. Note that
                  indices are 1-based an need to be converted for correct usage in
                  python.
                  
                  \textbf{Input:}
                  \begin{itemize}
                    \item \texttt{values numpy.ndarray}: 1D array of float64 values
                  \end{itemize}
                  \textbf{Output:}
                  \begin{itemize}
                    \item \texttt{np.array}: Permutation indices (1-based Fortran
                          indices)
                  \end{itemize}
                  \begin{lstlisting}[language=Python]
# - 1 at the end converts from fortrans 1-based indexing to pythons 0-based
indices = build_bst_index(values) - 1
                  \end{lstlisting}
                  
            \item \textbf{BST Range Query}
                  
                  Performs range queries on BST index. Accepts and returns Fortran 1-based
                  indices.
                  
                  \textbf{Input:}
                  \begin{itemize}
                    \item \texttt{values numpy.ndarray}: Original values array
                          
                    \item \texttt{indices numpy.ndarray}: BST index array from build\_bst\_index
                          (1-based)
                          
                    \item \texttt{lower\_bound float}: Lower bound of range (inclusive)
                          
                    \item \texttt{upper\_bound float}: Upper bound of range (inclusive)
                  \end{itemize}
                  \textbf{Output:}
                  \begin{itemize}
                    \item \texttt{Dictionary}: (matching\_indices, count) where indices
                          are 1-based
                  \end{itemize}
                  \begin{lstlisting}[language=Python]
# + 1 is added because of 0-based indexing. If the indices are already 1 based, this is not necessary
res = bst_range_query(values, indices + 1, lower_bound, upper_bound)
                  \end{lstlisting}
                  
            \item \textbf{Get Sorted Value}
                  
                  \textbf{Note:} This function is not implemented as a separate Python
                  function. Use direct array indexing instead:
                  \begin{lstlisting}[language=Python]
value = values[indices[position]]  # position is 1-based
                  \end{lstlisting}
          \end{enumerate}
          
    \item \textbf{Python k-d Tree Interface}
          
          \begin{enumerate}
            \item \textbf{Build k-d Tree Index}
                  
                  Python interface for constructing k-d trees. Requires Fortran-ordered
                  arrays and returns 1-based indices.
                  
                  \textbf{Input:}
                  \begin{itemize}
                    \item \texttt{points numpy.ndarray}: 2D array of points (dimensions
                          × n\_points) in Fortran order
                          
                    \item \texttt{dimension\_order numpy.ndarray}: Dimension ordering (1-based
                          Fortran indices, optional)
                  \end{itemize}
                  \textbf{Output:}
                  \begin{itemize}
                    \item \texttt{np.array}: Permutation indices (1-based Python indices)
                  \end{itemize}
                  \begin{lstlisting}[language=Python]
kd_indices = build_kd_index(points, dimension_order)
                  \end{lstlisting}
                  
            \item \textbf{Build Spherical k-d Tree}
                  
                  Alias for build\_kd\_index with the same parameters and behavior.
                  
                  \textbf{Input:} Same as build\_kd\_index \textbf{Output:} Same as build\_kd\_index
                  \begin{lstlisting}[language=Python]
sphere_indices = build_spherical_kd(vectors, dimension_order)
                  \end{lstlisting}
                  
            \item \textbf{Get Point from k-d Index}
                  
                  \textbf{Note:} This function is not implemented as a separate Python
                  function. Use direct array indexing instead:
                  \begin{lstlisting}[language=Python]
point = points[:, kd_indices[position - 1]]  # position is 1-based
                  \end{lstlisting}
          \end{enumerate}
\end{enumerate}